\begin{abstract}
I study the difficulty of evaluating small strategy changes in a live economic-game
competition through a deterministic, LLM-free agent. The agent placed sixth in the
GLEE agent track; its development fleet logged 264,432 games across bargaining,
negotiation and persuasion. In a simultaneous 8.6-hour comparison, five identical
policy builds accumulated approximately 3,000 games each yet differed by as much
as 88 overall rating points and 293 within a game family. This documents substantial
leaderboard dispersion without a policy change, while leaving the contributions of
initialization, matchmaking and a changing reference pool unresolved. Decision-level
replay and terminal acceptance counterfactuals help investigate the behavior behind
these scores: some interventions changed no actions, while transcript inspection
exposed message-parser and termination-limit failures. An observational comparison
with my second-place human-track play shows different participation behavior, but
does not isolate message-understanding ability. The results motivate simultaneous
controls, explicit rating-model assumptions and behavioral audits when evaluating
small interventions during short competition windows. Code and analysis scripts
are public; full empirical reproduction depends on the underlying game records.
\end{abstract}

\section{Introduction}
GLEE~\cite{shapira2024glee} evaluates strategic interaction in bargaining,
negotiation and persuasion. Its competition provides a setting in which to ask a
practical question: \emph{how can an entrant distinguish a useful policy change
from a noisy score change?} I address this through \textsc{7aidara\_Gamma}, a
deterministic agent that finished sixth with an overall rating of 2360.8, and
through my own human-track play, which finished second.

The contribution is an empirical case study of evaluation and behavior rather
than a new equilibrium solution. Automated-negotiation research already studies
agent evaluation, modular strategies and acceptance conditions~\cite{anac2011,baarslag2014decoupling,baarslag2013acceptance};
uncertainty in benchmark comparisons is also an established concern~\cite{agarwal2021}.
Here, five simultaneous identical-policy controls reveal rating dispersion in a
live competition, and offline audits distinguish changed behavior from changes
that never affected an action. Failed experiments and implementation errors are
part of the evidence, not merely development history.

\section{The agent and its operating assumptions}
A dispatcher routes games to three Python policies (858 / 660 / 781 lines),
configured by 114 tunables. Parameters are fixed during a run; within-game state
and belief updates are deterministic given the observed history. Seller
randomization is seeded from the game ID and round. A safety layer repairs
invalid actions and supplies a legal fallback after an exception. The reference
build is covered by 472 behavioral tests.\footnote{\textbf{LLM usage.} The competition
agent calls no language model. I used LLM coding assistance for analysis scripts
and for drafting and revising the paper.}

\textbf{Bargaining.} The policy uses the infinite-horizon proposer-share anchor
$(1-\delta_{\rm them})/(1-\delta_{\rm us}\delta_{\rm them})$~\cite{rubinstein1982},
with empirical floors, concession rules and explicit endgame handling. This is a
heuristic benchmark, not a best-response guarantee against arbitrary opponents.
Finite-horizon equilibrium predictions depend on the remaining horizon and seat;
the agent instead exploits the empirically favorable final-proposer position
(91.9\% agreement in a 1,415-game cohort), while accounting for termination risk.

\textbf{Negotiation.} Observed values came from $(80,100,120,150)$ scaled by
$\{1,10^2,10^4\}$, with both players on the same scale. The policy therefore
prices at possible tradable rungs. Independent uniform draws on that ladder
predict positive surplus in $6/16=37.5\%$ of value pairs; independence of roles
alone would not imply that rate. When observed information rules out profitable
trade, the policy walks away.

\textbf{Persuasion.} Consistent with GLEE's documented payoff structure, the
competition seller received the full price on a sale regardless of quality
(checked over 1,829 sold rounds). For a buyer who trusts a known recommendation
policy and purchases at posterior threshold $\tau$, recommending every high unit
and a fraction $q$ of low ones gives
$\Pr(H\mid\mathrm{recommend})=p/[p+(1-p)q]$ and
$q_{\max}=\min\{1,p(1-\tau)/[\tau(1-p)]\}$.
Because GLEE lacks commitment, this Bayesian-persuasion threshold~\cite{kamenica2011bayesian}
is a heuristic starting point, not a guarantee of optimality. The seller ramps its
budget with the horizon. The buyer combines message parsing with within-game
purchase outcomes and recommendation frequency; quality is hidden after a decline.

\section{Identical policies, different leaderboard histories}
Five byte-identical builds were polled together for 8.6 hours, producing 45
simultaneous snapshots over approximately 3,000 games per agent. The overall
min--max spread averaged 48.1 points and reached 88.0; the largest within-family
spread was 292.9 (Figure~\ref{fig:noise}). These are descriptive spreads, not
confidence intervals. The five agents had different initial ratings and match
histories, and successive polls are dependent. The experiment measures their
joint effect with sampling and pool changes, rather than isolating one source
of noise.

\begin{figure}[t]
\centering
\begin{tikzpicture}[x=0.95cm,y=0.95cm,line width=0.4pt]
\draw[gridline] (0.000,0.400) -- (5.550,0.400);
\node[axlab,anchor=east] at (-0.080,0.400) {2200};
\draw[gridline] (0.000,0.900) -- (5.550,0.900);
\node[axlab,anchor=east] at (-0.080,0.900) {2250};
\draw[gridline] (0.000,1.400) -- (5.550,1.400);
\node[axlab,anchor=east] at (-0.080,1.400) {2300};
\draw[gridline] (0.000,1.900) -- (5.550,1.900);
\node[axlab,anchor=east] at (-0.080,1.900) {2350};
\node[axlab,anchor=north] at (0.000,-0.04) {0h};
\node[axlab,anchor=north] at (1.285,-0.04) {2h};
\node[axlab,anchor=north] at (2.570,-0.04) {4h};
\node[axlab,anchor=north] at (3.855,-0.04) {6h};
\node[axlab,anchor=north] at (5.140,-0.04) {8h};
\draw[axline] (0.000,0) -- (5.550,0);
\draw[axline] (0.000,0) -- (0.000,2.300);
\fill[band] (0.000,1.027) -- (0.111,1.254) -- (0.225,1.164) -- (0.357,1.196) -- (0.490,1.261) -- (0.621,1.326) -- (0.740,1.430) -- (0.872,1.357) -- (1.004,1.346) -- (1.136,1.440) -- (1.268,1.381) -- (1.387,1.389) -- (1.505,1.370) -- (1.623,1.468) -- (1.755,1.491) -- (1.887,1.456) -- (2.005,1.464) -- (2.122,1.714) -- (2.255,1.612) -- (2.373,1.673) -- (2.506,1.626) -- (2.638,1.701) -- (2.770,2.045) -- (2.888,1.906) -- (3.021,1.882) -- (3.152,1.850) -- (3.285,1.860) -- (3.417,1.831) -- (3.536,1.883) -- (3.668,1.868) -- (3.786,1.930) -- (3.905,1.965) -- (4.023,1.907) -- (4.155,1.978) -- (4.274,2.144) -- (4.405,2.155) -- (4.537,2.092) -- (4.669,1.886) -- (4.787,1.863) -- (4.904,1.933) -- (5.036,1.929) -- (5.168,1.868) -- (5.286,1.862) -- (5.418,1.788) -- (5.550,1.843) -- (5.550,1.385) -- (5.418,1.293) -- (5.286,1.415) -- (5.168,1.464) -- (5.036,1.447) -- (4.904,1.513) -- (4.787,1.738) -- (4.669,1.684) -- (4.537,1.552) -- (4.405,1.601) -- (4.274,1.707) -- (4.155,1.714) -- (4.023,1.755) -- (3.905,1.626) -- (3.786,1.430) -- (3.668,1.451) -- (3.536,1.517) -- (3.417,1.511) -- (3.285,1.628) -- (3.152,1.460) -- (3.021,1.447) -- (2.888,1.296) -- (2.770,1.322) -- (2.638,1.346) -- (2.506,1.112) -- (2.373,1.090) -- (2.255,1.007) -- (2.122,0.963) -- (2.005,0.944) -- (1.887,1.008) -- (1.755,1.099) -- (1.623,1.035) -- (1.505,1.165) -- (1.387,1.042) -- (1.268,0.964) -- (1.136,0.929) -- (1.004,0.920) -- (0.872,0.774) -- (0.740,0.695) -- (0.621,0.599) -- (0.490,0.611) -- (0.357,0.504) -- (0.225,0.483) -- (0.111,0.375) -- (0.000,0.174) -- cycle;
\draw[c1] (0.000,0.438) -- (0.111,0.375) -- (0.225,0.687) -- (0.357,0.621) -- (0.490,0.662) -- (0.621,0.736) -- (0.740,0.883) -- (0.872,1.000) -- (1.004,0.973) -- (1.136,1.048) -- (1.268,1.275) -- (1.387,1.271) -- (1.505,1.253) -- (1.623,1.189) -- (1.755,1.349) -- (1.887,1.251) -- (2.005,1.282) -- (2.122,1.347) -- (2.255,1.201) -- (2.373,1.254) -- (2.506,1.112) -- (2.638,1.359) -- (2.770,1.322) -- (2.888,1.296) -- (3.021,1.447) -- (3.152,1.618) -- (3.285,1.711) -- (3.417,1.539) -- (3.536,1.588) -- (3.668,1.643) -- (3.786,1.547) -- (3.905,1.782) -- (4.023,1.796) -- (4.155,1.714) -- (4.274,1.707) -- (4.405,1.779) -- (4.537,1.802) -- (4.669,1.886) -- (4.787,1.738) -- (4.904,1.513) -- (5.036,1.447) -- (5.168,1.464) -- (5.286,1.415) -- (5.418,1.293) -- (5.550,1.385);
\draw[c2] (0.000,0.445) -- (0.111,0.778) -- (0.225,0.718) -- (0.357,0.819) -- (0.490,0.747) -- (0.621,0.807) -- (0.740,0.856) -- (0.872,0.774) -- (1.004,1.006) -- (1.136,0.929) -- (1.268,1.129) -- (1.387,1.042) -- (1.505,1.165) -- (1.623,1.035) -- (1.755,1.099) -- (1.887,1.008) -- (2.005,0.944) -- (2.122,0.963) -- (2.255,1.007) -- (2.373,1.090) -- (2.506,1.136) -- (2.638,1.451) -- (2.770,1.503) -- (2.888,1.618) -- (3.021,1.683) -- (3.152,1.850) -- (3.285,1.860) -- (3.417,1.817) -- (3.536,1.854) -- (3.668,1.782) -- (3.786,1.745) -- (3.905,1.965) -- (4.023,1.890) -- (4.155,1.729) -- (4.274,1.733) -- (4.405,1.601) -- (4.537,1.552) -- (4.669,1.684) -- (4.787,1.786) -- (4.904,1.741) -- (5.036,1.631) -- (5.168,1.653) -- (5.286,1.559) -- (5.418,1.615) -- (5.550,1.418);
\draw[c3] (0.000,1.027) -- (0.111,1.254) -- (0.225,1.164) -- (0.357,1.196) -- (0.490,1.261) -- (0.621,1.326) -- (0.740,1.430) -- (0.872,1.357) -- (1.004,1.346) -- (1.136,1.440) -- (1.268,1.381) -- (1.387,1.389) -- (1.505,1.370) -- (1.623,1.336) -- (1.755,1.221) -- (1.887,1.218) -- (2.005,1.096) -- (2.122,1.302) -- (2.255,1.317) -- (2.373,1.336) -- (2.506,1.312) -- (2.638,1.346) -- (2.770,1.599) -- (2.888,1.602) -- (3.021,1.586) -- (3.152,1.460) -- (3.285,1.673) -- (3.417,1.628) -- (3.536,1.517) -- (3.668,1.451) -- (3.786,1.430) -- (3.905,1.626) -- (4.023,1.755) -- (4.155,1.969) -- (4.274,2.144) -- (4.405,2.155) -- (4.537,2.092) -- (4.669,1.859) -- (4.787,1.782) -- (4.904,1.718) -- (5.036,1.838) -- (5.168,1.786) -- (5.286,1.599) -- (5.418,1.620) -- (5.550,1.415);
\draw[c4] (0.000,0.401) -- (0.111,0.646) -- (0.225,0.669) -- (0.357,0.654) -- (0.490,0.615) -- (0.621,0.599) -- (0.740,0.695) -- (0.872,0.923) -- (1.004,1.091) -- (1.136,1.207) -- (1.268,1.325) -- (1.387,1.346) -- (1.505,1.343) -- (1.623,1.468) -- (1.755,1.491) -- (1.887,1.456) -- (2.005,1.464) -- (2.122,1.714) -- (2.255,1.612) -- (2.373,1.478) -- (2.506,1.484) -- (2.638,1.436) -- (2.770,1.519) -- (2.888,1.567) -- (3.021,1.468) -- (3.152,1.597) -- (3.285,1.628) -- (3.417,1.511) -- (3.536,1.544) -- (3.668,1.600) -- (3.786,1.701) -- (3.905,1.829) -- (4.023,1.870) -- (4.155,1.978) -- (4.274,1.978) -- (4.405,2.008) -- (4.537,1.903) -- (4.669,1.802) -- (4.787,1.739) -- (4.904,1.933) -- (5.036,1.929) -- (5.168,1.868) -- (5.286,1.767) -- (5.418,1.569) -- (5.550,1.610);
\draw[c5] (0.000,0.174) -- (0.111,0.520) -- (0.225,0.483) -- (0.357,0.504) -- (0.490,0.611) -- (0.621,0.712) -- (0.740,0.865) -- (0.872,1.003) -- (1.004,0.920) -- (1.136,1.025) -- (1.268,0.964) -- (1.387,1.132) -- (1.505,1.208) -- (1.623,1.244) -- (1.755,1.200) -- (1.887,1.353) -- (2.005,1.346) -- (2.122,1.289) -- (2.255,1.601) -- (2.373,1.673) -- (2.506,1.626) -- (2.638,1.701) -- (2.770,2.045) -- (2.888,1.906) -- (3.021,1.882) -- (3.152,1.720) -- (3.285,1.829) -- (3.417,1.831) -- (3.536,1.883) -- (3.668,1.868) -- (3.786,1.930) -- (3.905,1.957) -- (4.023,1.907) -- (4.155,1.971) -- (4.274,2.082) -- (4.405,1.795) -- (4.537,1.766) -- (4.669,1.821) -- (4.787,1.863) -- (4.904,1.761) -- (5.036,1.835) -- (5.168,1.806) -- (5.286,1.862) -- (5.418,1.788) -- (5.550,1.843);
\node[axlab,rotate=90] at (-0.920,1.125) {Overall rating};
\node[ptlab,anchor=south west] at (0.050,2.270) {(a)};
\draw[gridline] (7.500,0.157) -- (13.050,0.157);
\node[axlab,anchor=east] at (7.420,0.157) {0};
\draw[gridline] (7.500,0.680) -- (13.050,0.680);
\node[axlab,anchor=east] at (7.420,0.680) {$+50$};
\draw[gridline] (7.500,1.203) -- (13.050,1.203);
\node[axlab,anchor=east] at (7.420,1.203) {$+100$};
\draw[gridline] (7.500,1.727) -- (13.050,1.727);
\node[axlab,anchor=east] at (7.420,1.727) {$+150$};
\node[axlab,anchor=north] at (7.500,-0.04) {0h};
\node[axlab,anchor=north] at (8.785,-0.04) {2h};
\node[axlab,anchor=north] at (10.070,-0.04) {4h};
\node[axlab,anchor=north] at (11.355,-0.04) {6h};
\node[axlab,anchor=north] at (12.640,-0.04) {8h};
\draw[axline] (7.500,0) -- (13.050,0);
\draw[axline] (7.500,0) -- (7.500,2.300);
\draw[c1] (7.500,0.157) -- (7.611,0.090) -- (7.725,0.418) -- (7.857,0.348) -- (7.990,0.391) -- (8.121,0.469) -- (8.240,0.623) -- (8.372,0.745) -- (8.504,0.717) -- (8.636,0.795) -- (8.768,1.033) -- (8.887,1.028) -- (9.005,1.009) -- (9.123,0.943) -- (9.255,1.111) -- (9.387,1.007) -- (9.505,1.040) -- (9.622,1.108) -- (9.755,0.956) -- (9.873,1.010) -- (10.006,0.862) -- (10.138,1.121) -- (10.270,1.082) -- (10.388,1.055) -- (10.521,1.213) -- (10.652,1.391) -- (10.785,1.489) -- (10.917,1.309) -- (11.036,1.360) -- (11.168,1.418) -- (11.286,1.317) -- (11.405,1.563) -- (11.523,1.578) -- (11.655,1.492) -- (11.774,1.484) -- (11.905,1.560) -- (12.037,1.585) -- (12.169,1.672) -- (12.287,1.517) -- (12.404,1.282) -- (12.536,1.213) -- (12.668,1.230) -- (12.786,1.179) -- (12.918,1.051) -- (13.050,1.148);
\draw[c2] (7.500,0.157) -- (7.611,0.505) -- (7.725,0.442) -- (7.857,0.548) -- (7.990,0.472) -- (8.121,0.535) -- (8.240,0.586) -- (8.372,0.501) -- (8.504,0.744) -- (8.636,0.663) -- (8.768,0.873) -- (8.887,0.781) -- (9.005,0.910) -- (9.123,0.774) -- (9.255,0.842) -- (9.387,0.746) -- (9.505,0.679) -- (9.622,0.698) -- (9.755,0.745) -- (9.873,0.832) -- (10.006,0.880) -- (10.138,1.210) -- (10.270,1.264) -- (10.388,1.384) -- (10.521,1.452) -- (10.652,1.627) -- (10.785,1.637) -- (10.917,1.593) -- (11.036,1.631) -- (11.168,1.556) -- (11.286,1.517) -- (11.405,1.748) -- (11.523,1.668) -- (11.655,1.501) -- (11.774,1.504) -- (11.905,1.367) -- (12.037,1.316) -- (12.169,1.453) -- (12.287,1.560) -- (12.404,1.513) -- (12.536,1.398) -- (12.668,1.420) -- (12.786,1.323) -- (12.918,1.381) -- (13.050,1.175);
\draw[c3] (7.500,0.157) -- (7.611,0.394) -- (7.725,0.300) -- (7.857,0.334) -- (7.990,0.402) -- (8.121,0.470) -- (8.240,0.579) -- (8.372,0.502) -- (8.504,0.491) -- (8.636,0.589) -- (8.768,0.527) -- (8.887,0.535) -- (9.005,0.515) -- (9.123,0.480) -- (9.255,0.360) -- (9.387,0.357) -- (9.505,0.229) -- (9.622,0.445) -- (9.755,0.461) -- (9.873,0.480) -- (10.006,0.455) -- (10.138,0.491) -- (10.270,0.755) -- (10.388,0.759) -- (10.521,0.741) -- (10.652,0.610) -- (10.785,0.832) -- (10.917,0.786) -- (11.036,0.669) -- (11.168,0.600) -- (11.286,0.579) -- (11.405,0.784) -- (11.523,0.919) -- (11.655,1.142) -- (11.774,1.326) -- (11.905,1.337) -- (12.037,1.271) -- (12.169,1.028) -- (12.287,0.947) -- (12.404,0.880) -- (12.536,1.006) -- (12.668,0.951) -- (12.786,0.755) -- (12.918,0.777) -- (13.050,0.563);
\draw[c4] (7.500,0.157) -- (7.611,0.413) -- (7.725,0.438) -- (7.857,0.422) -- (7.990,0.381) -- (8.121,0.364) -- (8.240,0.465) -- (8.372,0.703) -- (8.504,0.879) -- (8.636,1.001) -- (8.768,1.124) -- (8.887,1.146) -- (9.005,1.142) -- (9.123,1.274) -- (9.255,1.297) -- (9.387,1.261) -- (9.505,1.269) -- (9.622,1.531) -- (9.755,1.424) -- (9.873,1.284) -- (10.006,1.291) -- (10.138,1.240) -- (10.270,1.327) -- (10.388,1.378) -- (10.521,1.274) -- (10.652,1.409) -- (10.785,1.441) -- (10.917,1.319) -- (11.036,1.353) -- (11.168,1.412) -- (11.286,1.518) -- (11.405,1.651) -- (11.523,1.694) -- (11.655,1.807) -- (11.774,1.807) -- (11.905,1.839) -- (12.037,1.729) -- (12.169,1.624) -- (12.287,1.557) -- (12.404,1.760) -- (12.536,1.756) -- (12.668,1.692) -- (12.786,1.586) -- (12.918,1.379) -- (13.050,1.422);
\draw[c5] (7.500,0.157) -- (7.611,0.520) -- (7.725,0.481) -- (7.857,0.502) -- (7.990,0.615) -- (8.121,0.720) -- (8.240,0.881) -- (8.372,1.025) -- (8.504,0.938) -- (8.636,1.049) -- (8.768,0.984) -- (8.887,1.160) -- (9.005,1.240) -- (9.123,1.277) -- (9.255,1.231) -- (9.387,1.391) -- (9.505,1.384) -- (9.622,1.325) -- (9.755,1.651) -- (9.873,1.726) -- (10.006,1.678) -- (10.138,1.755) -- (10.270,2.115) -- (10.388,1.970) -- (10.521,1.945) -- (10.652,1.775) -- (10.785,1.890) -- (10.917,1.892) -- (11.036,1.946) -- (11.168,1.930) -- (11.286,1.996) -- (11.405,2.023) -- (11.523,1.971) -- (11.655,2.038) -- (11.774,2.154) -- (11.905,1.853) -- (12.037,1.823) -- (12.169,1.881) -- (12.287,1.925) -- (12.404,1.819) -- (12.536,1.895) -- (12.668,1.865) -- (12.786,1.924) -- (12.918,1.846) -- (13.050,1.904);
\node[axlab,rotate=90] at (6.380,1.125) {Change since start};
\node[ptlab,anchor=south west] at (7.550,2.270) {(b)};
\draw[c1,line width=0.8pt] (3.40,-0.62) -- (3.85,-0.62);
\node[axlab,anchor=west] at (3.92,-0.62) {B$_1$};
\draw[c2,line width=0.8pt] (4.75,-0.62) -- (5.20,-0.62);
\node[axlab,anchor=west] at (5.27,-0.62) {B$_2$};
\draw[c3,line width=0.8pt] (6.10,-0.62) -- (6.55,-0.62);
\node[axlab,anchor=west] at (6.62,-0.62) {B$_3$};
\draw[c4,line width=0.8pt] (7.45,-0.62) -- (7.90,-0.62);
\node[axlab,anchor=west] at (7.97,-0.62) {B$_4$};
\draw[c5,line width=0.8pt] (8.80,-0.62) -- (9.25,-0.62);
\node[axlab,anchor=west] at (9.32,-0.62) {B$_5$};
\end{tikzpicture}

\caption{Overall ratings of five identical policy builds over 8.6 hours.
Shading is the observed min--max envelope, not an uncertainty interval.
The mean spread is 48.1 points and its maximum is 88.0. Initial rating states
were not reset.}
\label{fig:noise}
\end{figure}

Competition scoring compares payoffs within a configuration and role and adjusts
for opponent strength. For interpreting late-stage updates, I used the simplified
EMA model $R_{t+1}=(1-\eta)R_t+\eta g_t$, with $\eta\approx0.002$ inferred from
rating deltas and $g_t$ denoting an adjusted game target. At constant $\eta$, its
half-life is $\log(0.5)/\log(1-\eta)\approx346$ family games. Under independent,
stationary targets, its effective sample size is $(2-\eta)/\eta=999$; lifetime
game volume therefore need not describe the information in the current rating.
These are model implications, not a reconstruction of every server update.

The reference pool changes as games enter it, making unmatched cross-time
comparisons confounded. Selecting or freezing a high-scoring slot also makes
stopping time part of the evaluation protocol. A further distinction concerns
the outcome being tested: for consecutive family updates,
\mbox{$\overline{\Delta R}=(R_{\rm end}-R_{\rm start})/n$}.
Mean rating change therefore tends toward zero as a rating settles, even for
a strong policy. An IID calculation based on such deltas cannot by itself
establish a sample requirement for detecting a difference in policy strength
(Appendix~\ref{app:measurement}).

\section{Behavioral measurements beyond the displayed rating}
\textbf{Terminal acceptance counterfactuals.} An available offer can be evaluated
exactly at its decision point: accepting banks
$\mathrm{share}\,\delta^{r-1}$. Comparing that value with the realized continuation
made 23,617 refusals informative about acceptance rules, including continuations
ending without agreement. This is a local, retrospective comparison; it does not
identify the value of every alternative policy or prove ex-ante optimality.
Repeated refusals within one game share an outcome.

\textbf{Decision-level replay.} The reference policy reproduced 4,272 of 4,272
bargaining and 13,056 of 13,056 negotiation decisions in the audited cohorts.
Replay then counted actions changed by a candidate flag, rather than how often
its condition held. One condition fired on 43\% of eligible games yet changed
0 of 13,056 decisions; another changed 0 of 13,479. Exact replay validates the
implementation on logged states, not the outcome of unobserved trajectories.

\textbf{Exploratory break-even fits.} Regressing rating delta on banked payoff
within a cell supplied diagnostic zero crossings (Appendix
Table~\ref{tab:breakeven}). Slopes were positive in 58 of 64 cells, with median
$R^2=0.60$ and split-half correlation $r=0.741$. These describe regularity in
the logged sample. A root depends on rating history, opponent adjustments and
the reference pool; it is not generally a recovered field median. A fourth
diagnostic estimates aggregate quality among declined persuasion rounds from
the known prior (Appendix~\ref{app:lim}).

\section{Interventions and implementation failures}
Across 30+ interventions and 49 runs, the clearest improvements concerned
implementation repairs and structural policy changes (Appendix
Table~\ref{tab:ledger}). Their evidence types differ; none of the empirical
comparisons alone establishes game-theoretic dominance. Seven consecutive tuning
arms yielded no reliable improvement in the comparisons used. A seller lie-rate
arm also showed no reliable difference across three runs ($t=-0.29$), which does
not establish equivalence.

Transcript inspection found purchases following messages such as ``Not this
one\ldots pass this round.'' The parser recognized ``not this round'' but not
``not this one'' and interpreted unmatched text as a recommendation. Re-scoring
123,700 buyer rounds reclassified 270 purchases, 258 of them low quality.
Another transcript exposed an acceptance gate opening at round 97 in negotiations
that terminated at 96. A separate message intervention changed take-up from
0.689 to 0.693; that observation does not establish that language generally has
no effect.

On the final morning, a fleet-wide concession change was followed by drops of
100--385 negotiation rating points over five hours. Without a concurrent
untreated control, the full decline cannot be attributed to the parameter alone.
Operationally, the lesson is to retain a control. Abandonment was also clustered:
80 of 109,881 rated games, with 78 in seven clock hours. Their recorded deltas
averaged $-6.76$ and summed to about $-541$; the sum is not their persistent
effect on an EMA rating.

\section{Agent Behavior Analysis}
\textbf{Observed behavior.} Bargaining settled at median round 2, with 1.5\%
no-deals and mean banked share 0.429. Negotiation walked away in 79.8\% of
untradable seats and agreed in 20--24\% of hidden-value games versus 85--89\%
with complete information. The persuasion buyer bought 17.4 units per game
where buying blind had positive expected value, versus 2.2 where it did not.
These are sample descriptions, not controlled treatment effects.

\textbf{Design alignment and failures.} The 472 tests encode intended actions
in named regimes. The reference snapshot passed them; the final-night working
tree, carrying experimental flags, failed 27. Together with replay, these tests
check implementation consistency rather than utility optimality. Four recurring
deviations were unreachable rules, a gate keyed to the wrong termination limit,
a permissive parser default and excessive last-round concession. In one bounded
negotiation cohort, 403 of 469 seller games reached a round-10 agreement at low
surplus (Appendix Table~\ref{tab:endgame}). Agreement rate alone concealed that
behavior because capitulation still counted as a deal.

\section{An observational same-author human baseline}
I played 2,690 rated human-track games and finished second. Sharing the game
families and opponent pool makes a useful case study, but does not match exact
opponents, messages or times. Table~\ref{tab:human} compares shared bargaining
cells and pooled persuasion regimes. Bargaining equally weights retained cells
matched on horizon type, seat and discount factors; persuasion pools by $mp$,
so within-regime configuration composition may still differ.

\begin{table}[t]
\centering\footnotesize
\caption{Observational human--agent comparison. Bargaining: mean banked share.
Persuasion: payoff per game in product-price units; $t$ has Welch's standard-error
form, signed human minus agent, without dependence adjustment. Persuasion rows
contain 48--92 human and 4,729--9,467 agent games.}
\label{tab:human}
\begin{tabular}{@{}lrrr@{\hspace{1.5em}}lrrr@{}}
\toprule
\multicolumn{4}{@{}c}{\textbf{Bargaining (known horizon)}} &
\multicolumn{4}{c@{}}{\textbf{Persuasion (buyer)}}\\
\cmidrule(r{1.5em}){1-4}\cmidrule(l){5-8}
seat & cells & human & agent & regime & human & agent & $t$\\
\midrule
player 1 & 14 & 0.444 & 0.447 & $mp>1$ & $+24.22$ & $+22.10$ & $+1.31$\\
player 2 & 13 & 0.512 & 0.720 & $mp=1$ & $+2.99$ & $+2.56$ & $+0.56$\\
 & & & & $mp<1$ & $-0.61$ & $+0.18$ & $-2.44$\\
\bottomrule
\end{tabular}
\end{table}

Among purchased items, the human's high-quality rate was 0.738 over 2,745
purchases and the agent's was 0.725 over 59,603. These rates concern selected
purchases, not message-reading accuracy. The agent's higher bargaining share in
the known-horizon second seat, and my negative mean persuasion payoff where
$mp<1$, are consistent with different participation and acceptance behavior.
They do not identify discipline rather than perception as a causal mechanism,
nor support a general claim about humans versus machines.

\section{Limitations and conclusion}
This is one account, one human participant and a changing competition pool.
The identical-policy experiment does not reset rating history or isolate
matchmaking; exploratory fits and pooled tests do not remove all dependence or
selection. Repeated tuning and flexible cell filters increase the risk of
selective findings. The public repository provides policies, tests and analysis
scripts, but full empirical reproduction requires private game records and
deployment metadata; anticipated organizer data do not by themselves resolve
that gap.

The useful finding is narrower than a universal failure of leaderboard
evaluation: substantial rating dispersion occurred without policy changes,
and behavioral audits exposed failures that score summaries alone obscured.
Simultaneous controls, explicit estimands, replay and transcript inspection
provide complementary evidence when live experimental windows are short.
\label{endmain}


